\answer{OA28}
总权重为 $\sum_{n=0}^\infty2^{-n}=2$，所以级数在每个 $x$ 处均收敛，且任意尾部的贡献都可统一控制。

先设 $x$ 为无理数。给定 $\varepsilon>0$，选 $N$ 使 $\sum_{n>N}2^{-n}<\varepsilon$。有限多个点 $r_0,\ldots,r_N$ 均不等于 $x$，故存在 $\delta>0$，使 $|y-x|<\delta$ 时这些点与 $y$ 的大小关系和与 $x$ 的大小关系相同。前 $N+1$ 项不变，于是
\[
 |f(y)-f(x)|\le\sum_{n>N}2^{-n}<\varepsilon.
\]
所以每个无理点都连续。

再设 $x=r_k$。分离第 $k$ 项，写成
\[
 f(y)=2^{-k}\mathbf1_{\{y>r_k\}}+g(y).
\]
对 $g$ 使用同样的有限前缀加尾部估计：前缀中已没有等于 $r_k$ 的点，故 $g$ 在 $r_k$ 连续。严格不等式使 $f(r_k)=g(r_k)$，从而
\[
 \boxed{\lim_{y\uparrow r_k}f(y)=f(r_k),\qquad
 \lim_{y\downarrow r_k}f(y)=f(r_k)+2^{-k}.}
\]
因此每个有理点都有大小为 $2^{-k}$ 的向上跳跃，左连续而不连续；全部连续点恰为无理数。跳点虽稠密，但尾部总权重很小，故不会破坏无理点的连续性。
